\section{Notations and basic facts}\label{section-basic}
We use the notations of Equation~\eqref{eq}, of Assumptions~\eqref{listintro1.1}-\eqref{listintro1.5} and of the introduction. In particular, we recall that $X=H^1_0(\Omega)\times L^2(\Omega)$ and
\[
D(A)=\left(H^2(\Omega)\cap H^1_0(\Omega)\right)\times H^1_0(\Omega)\,
A =
\begin{pmatrix}
0 & Id \\ \Delta-\alpha Id & -\gamma(x)
\end{pmatrix}\,.
\]
The operator $A$ is the classical linear damped wave operator corresponding to the linear part of~\eqref{eq}. Due to Lumer--Phillips theorem, we know that this operator generates a linear semigroup of contractions $e^{At}$ on $X$ and on $D(A)$ and that
\[
\forall t\geq 0\,,\,\left\|e^{At}\right\|_{\Lc(X)}\leq 1\quad\text{and}\quad\left\|e^{At}\right\|_{\Lc(D(A))}\leq 1\,.
\]
Notice that the second estimate is a direct consequence of the commutation of $A$ and $e^{At}$ and does not require any regularity on $\gamma$.

For any $\sigma\in [0,1]$, we set
\[
X^\sigma=\left(H^{1+\sigma}(\Omega)\cap H^1_0(\Omega)\right)\times
H^\sigma_0(\Omega)\,.
\]
Thus $X^0=X$ and $X^1=D(A)$ and $X^\sigma$ is an interpolation space between $X^0$ and $X^1$. In particular, by interpolation, $e^{At}$ is defined in $X^\sigma$ and we have
\begin{equation}\label{eq-bound-semigroup}
\forall \sigma\in [0,1]\,,\:\forall t\geq 0\,,\:
\left\|e^{At}\right\|_{\Lc(X^\sigma)}\leq 1\,.
\end{equation}

We set $F\in\Cc^0(X)$ to be the function
\begin{equation}\label{def-F}
F:U=\vect uv \in X\,\longmapsto\,\vect 0 {-f(\cdot,u)} \in X\,.
\end{equation}
Notice that, if $\Omega$ is two-dimensional, $H^1_0(\Omega) \hookrightarrow L^{p}(\Omega)$ for any $p\in[1,+\infty)$ and if $\Omega$ is three-dimensional $H^1_0(\Omega) \hookrightarrow L^6(\Omega)$. Thus, $f(u)$ is well defined in $L^2(\Omega)$ due to Assumption~\eqref{hyp-f-estim} if dim$(\Omega)=2$ or if dim$(\Omega)=3$ and $p\leq 3$. Moreover, for any $R$, $u$ and $v$ with $\|u\|_{H^1}\leq R$ and $\|v\|_{H^1}\leq R$, we have
\[
\|f(\cdot,u)-f(\cdot,v)\|_{L^2}=\left\|(u-v)\int_0^1
f'_u(\cdot,u+s(u-v))ds\right\|_{L^2}
\leq C(R)\|u-v\|_{H^1}
\]
and so $F$ is lipschitzian on the bounded sets of $X$. As a consequence, the damped wave equation~\eqref{eq} is well posed in $X$ and admits local solutions if dim$(\Omega)=2$ or if dim$(\Omega)=3$ and $p\leq 3$.

With the above notation, our main equation writes
\begin{equation}\label{eq2}
\partial_t U=AU + F(U)\quad U(t=0)=U_0\in X\,.
\end{equation}
In particular, Duhamel's formula yields
\[
U(t)=e^{At}U_0 + \int_0^t e^{A(t-s)}F(U(s))\,\drm s\,.
\]

We introduce the potential
\[
V(x,u)=\int_0^u f(x,s)ds\,.
\]
Due to~\eqref{hyp-f-estim} and the above arguments, $V(\cdot,u)$ defines a Lipschitz function from the bounded sets of $H^1_0(\Omega)$ into $L^1(\Omega)$. The classical energy associated to~\eqref{eq} is defined along a trajectory $U=(u,\partial_t u)$ as
\[
E(U)= \int_\Omega \frac12\left(|\grad u|^2+\alpha |u|^2+|\partial_t u|^2\right) +
V(x,u)\,.
\]
The damping effect appears by the computation
\begin{equation}\label{eq-deriv-E}
\partial_t E(U(t))= - \int_\Omega \gamma(x)|\partial_t u |^2\,.
\end{equation}
In particular, the energy $E$ is non-increasing along the trajectories. Moreover, the sign assumption~\eqref{hyp-f-sign} yields that $V(x,u)\geq 0$. Thus, we have that $E(U)\geq C \|U\|_X^2$ and that $E(t)$, $t\geq 0$, is bounded on the bounded sets of $X$. All together, the above properties show that for any $U_0\in X$, the solution $U=(u,\partial_t u)$ of~\eqref{eq} is defined for all non-negative times and remains in a bounded set of $X$, which only depends on $\|U_0\|_X$.

A fundamental question of this paper concerns the solution for which the energy is constant: are they equilibrium points or may they be moving trajectories? At least, the answer is known for the linear equation, see~\cite{Dafermos} and also~\cite{Haraux}.

\begin{theo}\label{th-stab-linear}
\emph{Dafermos (1978).}
Assume that the damping $\gamma\geq 0$ does not vanish everywhere. Then, for any $U_0\in X$, we have
\[
e^{At}U_0\xrightarrow[t\longrightarrow +\infty]{} 0\text{ in }X\,.
\]
\end{theo}



\section{Asymptotic compactness and reduction to a unique continuation problem}\label{section-asymp-compact}

In this section, we assume a fast enough semi-uniform linear decay as described by~\eqref{hyp-decay-1} and~\eqref{hyp-decay-2}. We first notice that, by linear interpolation, we have the following result.

\begin{prop}\label{prop-decay-interpole}
For any $\sigma_1,\sigma_2$ such that $0\leq \sigma_2<\sigma_1 \leq 1$, the linear semigroup is well defined from $X^{\sigma_1}$ in $X^{\sigma_2}$ and we have
\[
\forall U_0\in X^{\sigma_1}\,,\,\left\|e^{At}U_0\right\|_{X^{\sigma_2}} \leq
h(t)^{\sigma_1-\sigma_2} \|U_0\|_{X^{\sigma_1}}\,.
\]
\end{prop}
\begin{proof}
We interpolate the estimates~\eqref{eq-bound-semigroup} for $\sigma=1$ and~\eqref{hyp-decay-1} with respective weights $(\sigma_2/\sigma_1,1-\sigma_2/\sigma_1)$ and obtain
\[
\forall U_0\in D(A)\,,\,\left\|e^{At}U_0\right\|_{X^{\sigma_2/\sigma_1}} \leq
h(t)^{1-\sigma_2/\sigma_1}\|U_0\|_{D(A)}\,.
\]
It remains to interpolated the above estimate and~\eqref{eq-bound-semigroup} for $\sigma=0$ with respective weights $(\sigma_1,1-\sigma_1)$.
\end{proof}



We also need some regularity properties for $F$. The following properties depend on Sobolev embeddings and so of the dimension $d$ of $\Omega$. For $d=2$, which is the case in our examples, the properties are general. For $d=3$, they are more restrictive but they are shown in the same way. We choose to also consider this case in our paper for possible later uses.

\begin{prop}\label{prop-compact-F}
Assume that dim$(\Omega)=2$. Then for any $\sigma\in [0,1)$, the function $F$ maps any bounded set $\Bc$ of $X=H^1_0(\Omega)\times L^2(\Omega)$ in a bounded set $F(\Bc)$ contained in $X^\sigma= (H^{1+\sigma}(\Omega)\cap H^1_0(\Omega))\times H^\sigma_0(\Omega)$. Moreover, $F(\Bc)$ has compact enclosure in $X^\sigma$.

If dim$(\Omega)=3$ and if~\eqref{hyp-f-estim} holds for some $p\in [0,3)$, then the same properties hold for $\sigma\in [0,(3-p)/2)$.
\end{prop}
\begin{proof}
Assume that dim $(\Omega)=2$. First notice that we only have to show that $f(\cdot,u)$ is compactly bounded in $H^\sigma_0(\Omega)$ since the first component of $F$ is zero. Also notice that $f(x,0)=0$ due to the sign Assumption~\eqref{hyp-f-sign}, thus the Dirichlet boundary condition possibly contained in $H^\sigma_0(\Omega)$ will be fulfilled by $f(x,u)$ if $u\in H^1_0(\Omega)$. Due to the Sobolev embeddings, and since $\Omega$ is compact, it is sufficient to show that $F(\Bc)$ is bounded in $W^{1,q}(\Omega)$ for all $q\in [1,2)$ to obtain compactness in $H^\sigma(\Omega)$ for any $\sigma\in [0,1)$. Since $f$ is of polynomial type due to~\eqref{hyp-f-estim}, we know that $f(x,u)$ is bounded in $L^q(\Omega)$ for any $q\in [1,2)$. On the other hand, using~\eqref{hyp-f-estim}, we have
\begin{align*}
\|\grad(f(x,u))\|_{L^q} & \leq \|\grad_xf(x,u)\|_{L^q} + \left\|f'_u(x,u)\grad u\right\|_{L^q}\\
& \leq C \|(1+|u|)^p \|_{L^q} + C \left\|(1+|u|)^{p-1} \grad u\right\|_{L^q}\\
& \leq C\left(1+\|u\|_{L^{pq}}^p + \|\grad u\|_{L^2} \|u\|^{p-1}_{L^r}\right)
\end{align*}
with $r=(p-1)\frac {2q}{2-q}$ defined as soon as $q<2$. This shows that $f(\cdot,u)$ belongs to $W^{1,q}(\Omega)$ for any $q\in [1,2)$ and concludes the proof for dim$(\Omega)=2$.

The case dim$(\Omega)=3$ is similar once we use the suitable Sobolev embeddings.
\end{proof}


The main results of this section are the following asymptotic compactness properties.

\begin{prop}\label{prop-asymp-compact-1}
If dim$(\Omega)=2$, set $\sigma_*=\sigma_h$. If dim$(\Omega)=3$, assume that $p\in(1,3)$ in~\eqref{hyp-f-estim} and that $\sigma_h\in ((p-1)/2,1)$ in~\eqref{hyp-decay-2} and set $\sigma_*=\sigma_h - (p-1)/2$.

Let $U(t)=(u,\partial_t u)$ where $u$ solves~\eqref{eq} and let $(t_n)$ be a sequence of times such that $t_n\rightarrow +\infty$. Then, there exist a subsequence $(t_{\varphi(n)})$ and a solution $W(t)=(w,\partial_t w)(t)$ of~\eqref{eq} defined for all $t\in\RR$, such that
\[
\forall
t\in\RR\,,\,U(t_{\varphi(n)}+t)\,\xrightarrow[\:n\longrightarrow
+\infty\;]{}\,W(t)\;\text{ in }X=H^1_0(\Omega)\times L^2(\Omega)\,.
\]
Moreover, the solution $W$ is globally bounded in $X^\sigma$ for all $\sigma\in [0,\sigma_*)$ and the energy $E(W(t))$ is constant.
\end{prop}
\begin{proof}
Assume first that dim$(\Omega)=2$. We have
\begin{equation}\label{eq-prop-asymp-compact-1}
U(t_n)=e^{At_n}U_0 + \int_0^{t_n} e^{A(t_n-s)}F(U(s))\,\drm s\,.
\end{equation}
Due to Theorem~\ref{th-stab-linear}, the term $e^{At_n}U_0$ goes to zero in $X$. Thus, it remains to show that $\int_0^{t_n} e^{A(t_n-s)}F(U(s))\,\drm s$ is a compact term in $X$. First notice that $U(s)$ is uniformly bounded for $s\geq 0$ due to the non-increasing energy $E$ (see Section~\ref{section-basic}). Due to Proposition~\ref{prop-compact-F}, $F(U(s))$ thus belongs to a bounded set of $X^{\sigma_1}$ for all $\sigma_1\in [0,1)$. By Assumption~\eqref{hyp-decay-2} and Proposition~\ref{prop-decay-interpole}, $e^{A(t_n-s)}F(U(s))$ has an integral in $[0,t_n]$ bounded in $X^{\sigma_2}$ uniformly with respect to $n$, for any $\sigma_2 \in (0,\sigma_h)$. Thus $\int_0^{t_n} e^{A(t_n-s)}F(U(s))\,\drm s$ is a compact sequence in $X^\sigma$ for any $\sigma\in [0,\sigma_2)$. As a consequence, for any $\sigma$ as close as wanted to $\sigma_h$, we may extract a subsequence $(t_{\varphi(n)})$ such that $\int_0^{t_{\varphi(n)}} e^{A(t_{\varphi(n)}-s)}F(U(s))\,\drm s$ converges to some limit $W(0)$ in $X^\sigma$. Since the linear term of~\eqref{eq-prop-asymp-compact-1} goes to zero in $X$ for $t_{\varphi(n)}\rightarrow +\infty$, $U(t_{\varphi(n)})$ converges to $W(0)\in X^\sigma$ for the norm of $X$.

Let $W(t)=(w,\partial_t w)(t)$ be the maximal solution of the damped wave equation~\eqref{eq} corresponding to the initial data $W(0)$. Let $t\in\RR$, for $n$ large enough $t_{\varphi(n)}+t\geq 0$ and thus $U(t_{\varphi(n)}+t)$ is well defined and uniformly bounded in $X$. Since our equation is well posed, the solution is continuous with respect to the initial data. Thus, since $U(t_{\varphi(n)})$ converges to $W(0)$ in $X$, we have that $U(t_{\varphi(n)}+t)$ converges to $W(t)$ for all $t$ such that $W(t)$ is well defined. But due to the uniform bound on $U(t_{\varphi(n)}+t)$, $W(t)$ is uniformly bounded and thus the solution may be extended to a global solution $W(t)$, $t\in\RR$. In addition, the $X^\sigma$-bound obtained above for $W(0)$ only depends on the $X-$bound on $U(s)$ which is uniform due to non-increase of the energy of $U(t)$. Thus, the same arguments applied to the convergence $U(t_{\varphi(n)}+t)\rightarrow W(t)$ give the same $X^\sigma$-bound for $W(t)$ for all $t\in\RR$. Finally, since the energy of $U(t)$ is non-increasing and non-negative, for any $t\in\RR$, we must have $E(U(t_n+t))-E(U(t_n))\rightarrow 0$ when $n\rightarrow 0$ (since $t_n$ goes to $+\infty$). This shows that $E(W(t))$ is constant and finishes the proof.

The case dim$(\Omega)=3$ is similar once we take into account the constraints given by Proposition~\ref{prop-compact-F}.
\end{proof}


\goodbreak

\begin{prop}\label{prop-asymp-compact-2}
If dim$(\Omega)=2$, set $\sigma_*=\sigma_h$. If dim$(\Omega)=3$, assume that $p\in(1,3)$ in~\eqref{hyp-f-estim} and that $\sigma_h\in ((p-1)/2,1)$ in~\eqref{hyp-decay-2} and set $\sigma_*=\sigma_h - (p-1)/2$.

Let $\sigma>0$ and $R>0$. Let $U_n(t)=(u_n,\partial_t u_n)$ a sequence of solutions $u_n$ of~\eqref{eq} such that $(U_n(0))\subset X^\sigma$ and $\|U_n(0)\|_{X^\sigma}\leq R$. Let $(t_n)$ be a sequence of times such that $t_n\rightarrow +\infty$ and let $\sigma'\in [0,\min(\sigma,\sigma_*))$. Then, there exist subsequences $(t_{\varphi(n)})$ and $(U_{\varphi(n)})$ and a solution $W(t)=(w,\partial_t w)(t)$ of~\eqref{eq} defined for all $t\in\RR$, such that
\[
\forall
t\in\RR\,,\,U_{\varphi(n)}(t_{\varphi(n)}+t)\,\xrightarrow[\;n\longrightarrow
+\infty\:]{}_,W(t)\;\text{ in } X^{\sigma'}\,.
\]
Moreover, the solution $W$ is globally bounded in $X^{\sigma'}$ and the energy $E(W(t))$ is constant.
\end{prop}
\begin{proof}
The arguments are similar as the ones of the above proof of Proposition~\ref{prop-asymp-compact-1}. The term $e^{At_n}U_n(0)$ goes to zero in $X^{\sigma'}$ due to Proposition~\ref{prop-decay-interpole} because $\sigma'<\sigma$. We bound the integral $\int_0^{t_n} e^{A(t_n-s)}F(U_n(s))\,\drm s$ as in the proof of Proposition~\ref{prop-asymp-compact-1}: $U_n(s)$ is uniformly bounded in $X$, so $F(U_n(s))$ is uniformly bounded in $X^\eta$ with $\eta<1$ in dimension~$2$ or $\eta<(3-p)/2$ in dimension\,$3$. Proposition~\ref{prop-decay-interpole} together with~\eqref{hyp-decay-2} implies that the integral is uniformly bounded in $X^{\sigma'}$ with $\sigma'<\sigma_*$. The compactness follows by leaving any small amount of regularity in the process. To obtain the convergence to $W(t)$ for all $t$, we use the same argument as the one of the proof of Proposition~\ref{prop-asymp-compact-1} to first show the convergence in $X$. Then, the above arguments also show the compactness of $U(t_{\varphi(n)}+t)$ in $X^{\sigma'}$ and thus the convergence to $W(t)$ also holds in $X^{\sigma'}$. The last property is the same as the ones of Proposition~\ref{prop-asymp-compact-1}.
\end{proof}

The conclusions of Theorem~\ref{th1} then follow from Propositions~\ref{prop-asymp-compact-1} and~\ref{prop-asymp-compact-2} as soon as we can prove that $W\equiv 0$ for any subsequences of any sequences $(t_n)$ and $U_n$. To this end, notice that $E(W(t))$ is constant and its derivative~\eqref{eq-deriv-E} implies that $\int_\Omega \gamma(x) |\partial_t w|^2=0$ for all time. To formulate this property as a unique continuation property, we set as usual $z=\partial_t w$ and notice that $z$ solves
\begin{equation}\label{eq-UCP}
\begin{cases}
\partial_{tt}^2 z(x,t) = \Delta z(x,t) - \alpha z(x,t) - f'_u(x,w(x,t)) z\;&(x,t)\in\Omega\times\RR\\
z_{|\partial\Omega}(x,t)=0&(x,t)\in\partial\Omega\times \RR\\
\quad z(x,t)\equiv 0 & (x,t)\in \text{support}(\gamma)\times \RR
\end{cases}
\end{equation}
If this implies $z\equiv 0$ everywhere, this means that $w(x,t)=w(x)$ is constant in time and solves
\[
\Delta w(x)-\alpha w(x)=f(x,w(x))\,.
\]
Multiplying by $w$ and integrating, we obtain
\[
\|\grad w\|^2 +\alpha \|w\|^2 = -\int_\Omega f(x,w(x))w(x)\,\drm x\,.
\]
By the sign Assumption~\eqref{hyp-f-sign}, this yields $w\equiv 0$. Thus, it only remains to study this unique continuation property.

\begin{prop}\label{prop-UCP-implique-th}
Assume that $z\equiv 0$ is the only global solution of~\eqref{eq-UCP}. Then the decay assumptions~\eqref{hyp-decay-1} and~\eqref{hyp-decay-2} imply the conclusions of Theorem~\ref{th1}.
\end{prop}


